\usepackage[italian]{babel}

% Font attuale:
% - testo principale: STIX Two
% - sans-serif (titoli/elementi UI): Source Sans Pro
% Per tornare al font originale, sostituisci queste due righe con:
% \usepackage{libertinus}
% \usepackage[scaled=0.96]{sourcesanspro}
\usepackage{stix2}
\usepackage[scaled=0.96]{sourcesanspro}
\usepackage{lmodern}

\usepackage[style=italian/quotes]{csquotes}

\usepackage{booktabs} % tabelle
\usepackage{colortbl} % Per aggiungere colori alla tabella
\usepackage{xcolor}
\usepackage{graphicx}
\usepackage[most]{tcolorbox}

\graphicspath{{images/}}

\usepackage{float} % per posizionare le immagini esattamente dove indicator
\usepackage{scrhack} % evita un warning
\usepackage{etoolbox} % per compilazione condizionale v. righe seguenti
\usepackage{url}

\linespread{1.1}

\RedeclareSectionCommand[
  beforeskip=2.5\baselineskip,
  afterskip=1.5\baselineskip
]{chapter}

\RedeclareSectionCommand[
  beforeskip=1.5\baselineskip,
  afterskip=.8\baselineskip
]{section}

\setkomafont{chapter}{
    \sffamily
    \bfseries
    \color{RunBrandDark}
    \fontsize{28}{32}\selectfont
}

\setkomafont{section}{
    \sffamily
    \bfseries
    \color{RunBrand}
    \fontsize{16}{20}\selectfont
}

% elenchi puntati
\usepackage{enumitem}
\setlist{
  itemsep=0.4em,
  topsep=0.6em
}

\newbool{percorsi}
\boolfalse{percorsi}
% \booltrue{percorsi}

\newbool{covers}
% \boolfalse{covers}
\booltrue{covers}

\newbool{coverminimal}
% \booltrue{coverminimal}
\boolfalse{coverminimal}

% Preset font SOLO per copertina (prima/ultima pagina):
% - heros      : TeX Gyre Heros (più netto, meno arrotondato) [consigliato]
% - sourcesans : Source Sans Pro (attuale)
% - stix       : STIX Two Text (serif editoriale)
% - termes     : TeX Gyre Termes (serif classico)
\newcommand{\coverfontpreset}{heros}
\renewcommand{\coverfontpreset}{sourcesans}
% \renewcommand{\coverfontpreset}{stix}
% \renewcommand{\coverfontpreset}{termes}

% Variante tabella progressione:
% - attuale    : tabella boxed con badge settimana
% - precedente : tabella colorata semplice
% - nuova      : tabella editoriale con colonna "fase"
\newcommand{\progressiontablevariant}{attuale}
% \renewcommand{\progressiontablevariant}{precedente}
% \renewcommand{\progressiontablevariant}{nuova}

% Variante indice:
% - capitoli   : mostra solo i capitoli (consigliato per opuscolo)
% - dettagliato: mostra capitoli + sezioni
\newcommand{\tocvariant}{capitoli}
% \renewcommand{\tocvariant}{dettagliato}

\newcommand{\applytocdepth}{
    \ifdefstring{\tocvariant}{dettagliato}{
        \setcounter{tocdepth}{1}
    }{
        \setcounter{tocdepth}{0}
    }
}

\newcommand{\coverfontsourcesans}{\sffamily}
\newfontfamily\coverfontheros{TeX Gyre Heros}[Scale=MatchLowercase]
\newcommand{\coverfontstix}{\rmfamily}
\newfontfamily\coverfonttermes{TeX Gyre Termes}[Scale=MatchLowercase]

\newcommand{\coverfonttitle}{\coverfontheros}
\newcommand{\coverfonttext}{\coverfontheros}

\ifdefstring{\coverfontpreset}{sourcesans}{
    \renewcommand{\coverfonttitle}{\coverfontsourcesans}
    \renewcommand{\coverfonttext}{\coverfontsourcesans}
}{}
\ifdefstring{\coverfontpreset}{stix}{
    \renewcommand{\coverfonttitle}{\coverfontstix}
    \renewcommand{\coverfonttext}{\coverfontstix}
}{}
\ifdefstring{\coverfontpreset}{termes}{
    \renewcommand{\coverfonttitle}{\coverfonttermes}
    \renewcommand{\coverfonttext}{\coverfonttermes}
}{}

\usepackage{verbatim} % for comments
% \usepackage{listings} % for comments
\usepackage{xcolor} % for colorbox

\usepackage{url} % for URL

\ifbool{covers}{
\usepackage{tikz}
\usepackage{relsize} % for \smaller
}

\usepackage{hyperref} % per i link nel sommario

%\setlength{\parindent}{10pt}
%\setlength{\parskip}{1.4ex plus 0.35ex minus 0.3ex}
%\setlength{\parskip}{1.4ex plus 0.35ex minus 0.3ex}

\newcommand{\q}[1]{>>\textit{#1}<<}

% ===== Style system (veste grafica) =====
\definecolor{RunBrand}{HTML}{E4572E}
\definecolor{RunBrandDark}{HTML}{B23A1B}
\definecolor{RunInk}{HTML}{1F2937}
\definecolor{RunSoft}{HTML}{F8FAFC}
\definecolor{RunTipBg}{HTML}{FFF7ED}
\definecolor{RunWarnBg}{HTML}{FEF2F2}
\definecolor{RunNextBg}{HTML}{ECFDF5}

\hypersetup{
    colorlinks=true,
    linkcolor=RunBrandDark,
    urlcolor=RunBrand,
    citecolor=RunBrandDark
}

\setkomafont{disposition}{\sffamily\bfseries\color{RunBrandDark}}
\addtokomafont{pageheadfoot}{\color{RunInk}}

% Stile indice
\addto\captionsitalian{\renewcommand{\contentsname}{Indice}}
\setkomafont{chapterentry}{\sffamily\bfseries\color{RunBrandDark}}
\setkomafont{chapterentrypagenumber}{\sffamily\bfseries\color{RunBrandDark}}

\newtcolorbox{chapterleadbox}{
    enhanced,
    colback=RunSoft,
    colframe=RunBrand,
    boxrule=0.7pt,
    arc=1.2mm,
    left=1.2mm,
    right=1.2mm,
    top=1.2mm,
    bottom=1.2mm,
    before skip=0.6\baselineskip,
    after skip=1.1\baselineskip
}

% Box callout con titolo dinamico:
% uso interno, normalmente si usano le macro \consiglio, \errorecomune, \prossimopasso.
\newtcolorbox{tipbox}[1]{
    enhanced,
    colback=RunTipBg,
    colframe=RunBrand,
    title={\sffamily\bfseries #1},
    fonttitle=\small,
    boxrule=0.6pt,
    arc=1mm,
    left=1.1mm,
    right=1.1mm,
    top=1mm,
    bottom=1mm,
    before skip=0.8\baselineskip,
    after skip=0.8\baselineskip
}

% Box callout per avvisi/criticità:
% uso interno, normalmente si usa la macro \errorecomune.
\newtcolorbox{warnbox}[1]{
    enhanced,
    colback=RunWarnBg,
    colframe=RunBrandDark,
    title={\sffamily\bfseries #1},
    fonttitle=\small,
    boxrule=0.6pt,
    arc=1mm,
    left=1.1mm,
    right=1.1mm,
    top=1mm,
    bottom=1mm,
    before skip=0.8\baselineskip,
    after skip=0.8\baselineskip
}

% Box callout per azioni concrete e CTA:
% uso interno, normalmente si usano \prossimopasso e \funnelcta.
\newtcolorbox{nextstepbox}[1]{
    enhanced,
    colback=RunNextBg,
    colframe=RunBrandDark,
    title={\sffamily\bfseries #1},
    fonttitle=\small,
    boxrule=0.6pt,
    arc=1mm,
    left=1.1mm,
    right=1.1mm,
    top=1mm,
    bottom=1mm,
    before skip=0.8\baselineskip,
    after skip=0.8\baselineskip
}

% Apertura capitolo con blocco "cosa otterrai".
% Esempio:
% \chapterlead{In questo capitolo scoprirai:\begin{itemize}\item ...\end{itemize}}
\newcommand{\chapterlead}[1]{
    \begin{chapterleadbox}
        #1
    \end{chapterleadbox}
}

% Evidenzia un suggerimento pratico.
% Esempio: \consiglio{Inizia con 20 minuti a giorni alterni.}
\newcommand{\consiglio}[1]{
    \begin{tipbox}{Consiglio pratico}
        #1
    \end{tipbox}
}

% Evidenzia un errore da evitare.
% Esempio: \errorecomune{Partire troppo forte alla prima settimana.}
\newcommand{\errorecomune}[1]{
    \begin{warnbox}{Errore comune}
        #1
    \end{warnbox}
}

% Chiude una sezione/capitolo con un'azione immediata.
% Esempio: \prossimopasso{Nella prossima uscita alterna 1' corsa + 4' cammino.}
\newcommand{\prossimopasso}[1]{
    \begin{nextstepbox}{Prossimo passo}
        #1
    \end{nextstepbox}
}

% CTA funnel di fine capitolo (titolo + testo/URL/QR).
% Esempio:
% \funnelcta{Scarica il piano 4 settimane}{Vai su: tuosito.it/piano}
\newcommand{\funnelcta}[2]{
    \begin{nextstepbox}{Continua il percorso}
        {\bfseries\sffamily #1}\\[0.2em]
        #2
    \end{nextstepbox}
}

\newcommand{\weekbadge}[1]{
    \tcbox[
        on line,
        boxsep=1pt,
        left=4pt,
        right=4pt,
        top=2pt,
        bottom=2pt,
        arc=1.5mm,
        colback=RunBrandDark,
        colframe=RunBrandDark
    ]{\textcolor{white}{\sffamily\bfseries #1}}
}

\newcommand{\focusbadge}[1]{
    \tcbox[
        on line,
        boxsep=1pt,
        left=3pt,
        right=3pt,
        top=1.5pt,
        bottom=1.5pt,
        arc=1.2mm,
        colback=RunBrand!10,
        colframe=RunBrand!35
    ]{\textcolor{RunBrandDark}{\sffamily\bfseries #1}}
}

\newcommand{\walkword}{\textcolor{RunInk}{camminata}}
\newcommand{\runword}{\textcolor{RunBrandDark}{corsa}}

\newcommand{\progressiontableprevious}{
    \small
    \setlength{\tabcolsep}{6pt}
    \renewcommand{\arraystretch}{1.25}
    \rowcolors{2}{RunBrand!8}{white}
    \begin{tabular}{>{\centering\arraybackslash}p{0.18\linewidth} >{\raggedright\arraybackslash}p{0.66\linewidth}}
    \toprule
    \rowcolor{RunBrandDark}
    \textcolor{white}{\sffamily\bfseries Sett.} & \textcolor{white}{\sffamily\bfseries Allenamento} \\
    \midrule
    1  & 10 volte: 5' \walkword{} + 1' \runword{} \\
    2  & 10 volte: 4' \walkword{} + 2' \runword{} \\
    3  & 10 volte: 3' \walkword{} + 3' \runword{} \\
    4  & 6 volte: 5' \walkword{} + 5' \runword{} \\
    5  & 3 volte: 10' \walkword{} + 10' \runword{} \\
    6  & 2 volte: 10' \walkword{} + 20' \runword{} \\
    7  & 5' \walkword{} + 25' \runword{} \\
    8  & 5' \walkword{} + 35' \runword{} \\
    9  & 5' \walkword{} + 45' \runword{} \\
    10 & 1 ora di \runword{} \\
    \bottomrule
    \end{tabular}
}

\newcommand{\progressiontablecurrent}{
    \begin{tcolorbox}[
        enhanced,
        width=0.96\linewidth,
        colback=RunSoft,
        colframe=RunBrand,
        boxrule=0.7pt,
        arc=1.5mm,
        left=1.5mm,
        right=1.5mm,
        top=1.5mm,
        bottom=1.5mm
    ]
    \small
    \setlength{\tabcolsep}{6pt}
    \renewcommand{\arraystretch}{1.32}
    \rowcolors{2}{RunBrand!8}{white}
    \begin{tabular}{>{\centering\arraybackslash}p{0.20\linewidth} >{\raggedright\arraybackslash}p{0.64\linewidth}}
    \toprule
    \rowcolor{RunBrandDark}
    \textcolor{white}{\sffamily\bfseries Sett.} & \textcolor{white}{\sffamily\bfseries Allenamento} \\
    \midrule
    \weekbadge{1}  & 10 volte: 5' \walkword{} + 1' \runword{} \\
    \weekbadge{2}  & 10 volte: 4' \walkword{} + 2' \runword{} \\
    \weekbadge{3}  & 10 volte: 3' \walkword{} + 3' \runword{} \\
    \weekbadge{4}  & 6 volte: 5' \walkword{} + 5' \runword{} \\
    \weekbadge{5}  & 3 volte: 10' \walkword{} + 10' \runword{} \\
    \weekbadge{6}  & 2 volte: 10' \walkword{} + 20' \runword{} \\
    \weekbadge{7}  & 5' \walkword{} + 25' \runword{} \\
    \weekbadge{8}  & 5' \walkword{} + 35' \runword{} \\
    \weekbadge{9}  & 5' \walkword{} + 45' \runword{} \\
    \weekbadge{10} & 1 ora di \runword{} \\
    \bottomrule
    \end{tabular}
    \end{tcolorbox}
}

\newcommand{\progressiontablenew}{
    \begin{tcolorbox}[
        enhanced,
        width=0.96\linewidth,
        colback=white,
        colframe=RunBrandDark,
        boxrule=0.9pt,
        arc=2mm,
        left=1.2mm,
        right=1.2mm,
        top=1.2mm,
        bottom=1.2mm
    ]
    {\sffamily\bfseries\color{RunBrandDark} Progressione orientativa · 10 settimane}\par\vspace{0.45em}
    \footnotesize
    \setlength{\tabcolsep}{5pt}
    \renewcommand{\arraystretch}{1.24}
    \rowcolors{2}{RunSoft}{RunBrand!5}
    \begin{tabular}{>{\centering\arraybackslash}p{0.15\linewidth} >{\centering\arraybackslash}p{0.22\linewidth} >{\raggedright\arraybackslash}p{0.47\linewidth}}
    \rowcolor{RunBrandDark}
    \textcolor{white}{\sffamily\bfseries Sett.} & \textcolor{white}{\sffamily\bfseries Fase} & \textcolor{white}{\sffamily\bfseries Allenamento} \\
    \weekbadge{1}  & \focusbadge{Avvio}      & 10 volte: 5' \walkword{} + 1' \runword{} \\
    \weekbadge{2}  & \focusbadge{Ritmo}      & 10 volte: 4' \walkword{} + 2' \runword{} \\
    \weekbadge{3}  & \focusbadge{Equilibrio} & 10 volte: 3' \walkword{} + 3' \runword{} \\
    \weekbadge{4}  & \focusbadge{Svolta}     & 6 volte: 5' \walkword{} + 5' \runword{} \\
    \weekbadge{5}  & \focusbadge{Base}       & 3 volte: 10' \walkword{} + 10' \runword{} \\
    \weekbadge{6}  & \focusbadge{Tenuta}     & 2 volte: 10' \walkword{} + 20' \runword{} \\
    \weekbadge{7}  & \focusbadge{Continuità} & 5' \walkword{} + 25' \runword{} \\
    \weekbadge{8}  & \focusbadge{Autonomia}  & 5' \walkword{} + 35' \runword{} \\
    \weekbadge{9}  & \focusbadge{Fiducia}    & 5' \walkword{} + 45' \runword{} \\
    \weekbadge{10} & \focusbadge{Traguardo}  & 1 ora di \runword{} \\
    \end{tabular}
    \end{tcolorbox}
}

\newcommand{\renderprogressiontable}{
    \ifdefstring{\progressiontablevariant}{precedente}{
        \progressiontableprevious
    }{
        \ifdefstring{\progressiontablevariant}{nuova}{
            \progressiontablenew
        }{
            \progressiontablecurrent
        }
    }
}

\title{Fatti un regalo: inizia a correre}
\author{Daniele Pallastrelli} 
\date{\today} 

